\chapter{Start Here}

This is a self-contained path toward broad Latin reading, with independent
reading of the 1962 \emph{Missale Romanum} as its principal destination. The
course begins with ordinary project-composed sentences about visible things
and actions, then widens into Scripture, Cicero, Augustine, other excellent
classical and Christian prose, theology, canon law, poetry, and the varied
registers of the Missal. Every assigned excerpt, explanation, vocabulary set,
memory sheet, exercise, answer, and model belongs to the course bundle.

The common trunk contains twenty-two substantial packets. They group the
forty-six established lessons without renumbering their stable lesson IDs.
After the common advanced trunk, eight separate poetry packets form a branch
for deeper work in prosody, classical meters, Christian and rhythmic verse,
commentary, imitation, and composition. The trunk itself already introduces
poetic syntax and supplies the poetry needed for reading the Missal.

Begin with this guide, keep the Reference Grammar and Reading Lexicon nearby,
and use the Memory Workbook whenever fuller detachable repetition is useful.
The four cumulative self-review companions provide fresh work at broad
turning points. Use them when a fresh view of accumulated abilities would make
the next direction clearer.

\section{How every trunk packet flows}

After a compact cover and contents, every packet proceeds in the same order:

\begin{enumerate}
  \item \textbf{Teaching}: two or three connected stable lessons, each with
        definitions, full explanation, paradigms or decision procedures,
        worked examples, contrasts, and common errors;
  \item \textbf{Memory Work}: populated new, short-interval, longer-interval,
        and cumulative vocabulary and grammar retrieval;
  \item \textbf{Reading and Composition}: ordinary and source Latin chosen for
        the available grammar, followed by a cumulative step in writing Latin;
  \item \textbf{Cumulative Practice}: all four authoritative variants for
        every stable lesson, followed by a qualitative description of what
        should be reasonably dependable before continuing;
  \item \textbf{Answers and Models}: answers drawn from the same sources as the
        memory work, readings, compositions, and worksheets; and
  \item \textbf{Scope and Sources}: the packet boundary, references actually
        used, and generation record.
\end{enumerate}

The cover identifies the packet key, its stable lessons, what it assumes, what
it develops, and what follows. Later packets use terminology and vocabulary
already taught; the reference companions provide a retrieval home without
replacing the packet's own teaching.

\section{A study loop built on repetition}

\StudyLoop

There are no weeks, quotas, or fixed session lengths. A transparent section may
move quickly; a dense paradigm or author may deserve several returns. Use the
four variants and the Memory Workbook as repetition, not as bureaucracy. A
form or word is becoming secure when it can be retrieved in a fresh context,
its decisive rule can be explained, and an error can be noticed and repaired.

Keep a simple error record under five headings: forms; agreement and
government; clause structure; vocabulary and idiom; composition and style.
For a substantive correction, record the first decision, the corrected
analysis, the evidence that decides between them, and one new example. The
label “careless” explains nothing; a named rule and a fresh example make the
correction reusable.

\section{Vocabulary is a cumulative course strand}

Learn a noun with its dictionary form and gender, a verb with its principal
parts and constructions, and other words with the information needed to
recognize and use them. A single English gloss is not enough. The populated
memory rounds ask for recognition in both directions, dictionary-form
reconstruction, paradigms, word families, semantic contrasts, collocations,
governed cases, and use in connected Latin.

New vocabulary returns in examples, the next packet, later retrieval rounds,
and cumulative readings. Use the compact packet sheets first; use the Memory
Workbook for full detachable rounds, extra blank copies, and mixed review.
Repetition is doing its work when recall becomes prompt enough that attention
can remain on the sentence and its thought.

\section{Reading before lookup}

Before consulting a reference, mark the finite verb or verbal center, propose
dictionary forms, group words by agreement and government, and record more
than one formally possible reading where an ending is ambiguous. Then look up
only what remains uncertain. A familiar English rendering does not replace a
parse, and a parser cannot make the grammatical decision for you.

For unfamiliar text, distinguish words known directly, words inferable from
form or family, and words that genuinely need lookup. Record the hypothesis
before the lookup. The aim is not zero dictionary use; it is accurate,
purposeful lookup that no longer substitutes for reading.

\section{Composition begins at the beginning}

Early writing changes number, person, tense, or one governed word; it then
builds phrases and transparent sentences about ordinary subjects. Later work
combines clauses, transforms constructions, translates under constraints, and
imitates identified models. The advanced trunk develops paragraph structure,
periodic prose, register, documentation, and repeated revision before asking
for independent composition and research. The poetry branch applies the same
progression to verse only after its opening assumptions are reasonably secure.

\section{Before continuing}

Each packet closes by naming the abilities to make reasonably dependable and
the exact work worth repeating. Use fresh Latin and ask whether you can
reliably recognize the relevant forms, explain the evidence for a decision,
read connected sentences without leaning
on remembered English, and use the same patterns in your own Latin.
Occasional slips that you can detect and repair differ from repeated
uncertainty about the governing rule. When uncertainty recurs, return to the
exact explanation or memory sheet named by the packet and then try different
Latin. Continue when the listed abilities are reasonably dependable.

\section{How to use answers}

Keep the \emph{Answers and Models} part closed during a first attempt. For a
closed item, find the form, function, government, clause relation, or source
evidence that decides the answer; a matching paraphrase alone is not enough.
For open composition, compare morphology, syntax, sense, cohesion, idiom, and
register. A printed model is not the only possible good Latin.

Some translations admit several English wordings. An alternative is sound
when it preserves agency, negation, time, dependency, and material ambiguity.
Record and defend such an alternative instead of changing it merely to match
the model's style.

\section{Source classes and course limits}

Received passages carry exact citations or stable IDs mapped to the passage
inventory. An adaptation is labeled as an adaptation. Uncited short drills and
ordinary examples are project-composed unless a local note says otherwise.
English translations and glosses are project translations unless an
identified human translation is named. Never treat simplified wording as the
exact prose of Cicero, Augustine, Scripture, the Missal, or another source.

The course teaches reading, grammar, vocabulary, style, and composition, not
pronunciation. Quoted Missal text follows the identified 1962 edition and its
recorded normalization. Other authors may use different ligatures,
capitalization, punctuation, consonantal \Latin{i/j}, or \Latin{u/v}
conventions. Copy an assigned source as directed, but do not let typography
decide a grammatical function.
